% Created 2024-03-28 Thu 18:03
% Intended LaTeX compiler: pdflatex
\documentclass[11pt]{article}
\usepackage[utf8]{inputenc}
\usepackage[T1]{fontenc}
\usepackage{graphicx}
\usepackage{longtable}
\usepackage{wrapfig}
\usepackage{rotating}
\usepackage[normalem]{ulem}
\usepackage{amsmath}
\usepackage{amssymb}
\usepackage{capt-of}
\usepackage{hyperref}
\usepackage{booktabs}
\usepackage{minted}
\author{James Brusey and Elena Gaura}
\date{\today}
\title{AI and ML in Soil Agriculture}
\hypersetup{
 pdfauthor={James Brusey and Elena Gaura},
 pdftitle={AI and ML in Soil Agriculture},
 pdfkeywords={},
 pdfsubject={},
 pdfcreator={Emacs 29.2 (Org mode 9.5.5)}, 
 pdflang={English}}
\begin{document}

\maketitle
\bibliography{/Users/james/Documents/zotero-export}

\section{Challenge to feed 10 billion by 2050}
\label{sec:org8a573b6}
\begin{itemize}
\item \cite{searchingerCreatingSustainableFood2019} suggests that this goal is possible but there are many specific challenge areas, such as reducing food waste, and increasing production without increasing land use.
\item Aside from reforesting unproductive agricultural areas, we also need to protect and restore peatlands.
\item We also need to reduce the damaging affect from pesticide and fertiliser run-off into natural eco-systems.
\item 
\end{itemize}
\section{Soils are crucial for future food production}
\label{sec:org5319097}
\begin{itemize}
\item Over-watering can deplete soils of nutrients and raise the salt-table.
\item 
\end{itemize}
\section{AI and ML to the rescue?}
\label{sec:org3f851ef}
\begin{itemize}
\item Not likely to replace existing methods but maybe complementary
\item 
\end{itemize}
\section{AI for measuring soil health}
\label{sec:org9c56d30}
\begin{itemize}
\item Image analysis can pick up variations in crops to enable underlying problems with soil health to be exposed.
\item 
\end{itemize}
\section{Trends for AI in agriculture}
\label{sec:org4608252}
\begin{itemize}
\item Initial work in solving specific problems - such as killing weeds with lasers.
\item Holistic approaches to whole community solutions. Need to make it work for the community - not focus on techno-solutions.
\item Google maps parallel - need solution that is equivalent - makes every farmer small percentage more efficient.
\item As opposed to more techno solutions - lasering weeds - which may only be viable in limited situations.
\end{itemize}


\section{Points}
\label{sec:orgc9cba20}
\begin{itemize}
\item Solutions for robotic pickers and crop yield assessment
\item 
\end{itemize}
\end{document}